% sec-07: Document 07, the federal compliance guide for the LLM and robotic
% workflow.  This is the operative document for master prompt clause C.
\docpart{FDA Compliance Guide}{FDA LLM and Robotic Workflow National Compliance
Guide for a Phase 1 PDAC Clinical Trial Site}

\section{Purpose and Scope}

\draftnote{Say what this document is and what it is not. It is a map from the
site to the federal requirements that already exist. It is not a submission, not
a guidance document, and not an FDA position. State that in the first paragraph,
because the master prompt's clause C asks for the workflow to be acceptable to
the agency and the honest form of that ask is a map, not a claim.}

\section{The Regulatory Position of the Site}

\draftnote{Take the combined investigational new drug and investigational device
exemption position from \mvfile{trial-protocol/} and
\mvfile{trial-ind/}, and the filing itself from paper 16 of
\mvfile{funding/move-in/inputs/READMES/README-LLM-Pancreatic-Oncology-Clinical-Trial-System-Large-Documents-Funding-and-AI-Peer-Review.md},
the 172-page Investigational New Drug Application. State which parts of the
workflow sit under the drug filing, which under the device pathway, and which
under neither because they are practice of medicine.}

\section{Correction Map to 21 CFR Parts 11, 50, 54, 312, and 812}

\begin{mvtable}
\begin{tabularx}{\textwidth}{>{\raggedright\arraybackslash}p{2.4cm} >{\raggedright\arraybackslash}p{4.2cm} >{\raggedright\arraybackslash}X}
\headrow \thead{Citation} & \thead{Requirement} & \thead{Where this site discharges it}\\
\midrule
21 CFR 11.10 & to be fixed & to be fixed \\
21 CFR 312.62 & to be fixed & to be fixed \\
\end{tabularx}
\end{mvtable}
\tabcapl{tab:fdamap}{The correction map. Stage 2 completes it from the three
adapted frameworks at \mvfile{regulatory/}, one row per requirement, and every
row names a document in this package.}

\draftnote{Build the map from
\mvfile{regulatory/adaption-ich-e6r3/}, \mvfile{regulatory/Adaption-21-CFR-Part-50/}
and \mvfile{regulatory/Adaption-21-CFR-Part-312/}. Add 21 CFR Part 11 for the
electronic record and signature position and 21 CFR Part 54 for the financial
disclosure position, which is the one that decides the sponsor-investigator
firewall in document 14 \sect{}3.}

\section{ICH E6(R3) Mapping}

\draftnote{Map the site to the good clinical practice principles, using the
author's own adaptation rather than the source guideline where the two differ,
and say when they differ.}

\section{The Advisory-Only Boundary and Software Change Control}

\draftnote{This is the load-bearing section of the document. State the boundary
as a rule that can be tested: what the model may produce, what a human must do
before any of it reaches a participant, and what record proves the human acted.
Take the verification-before-generation position from paper 11 and the
supervision rationale from paper 13 of the deck table. State the change control
plan against the FDA predetermined change control plan guidance.}

\section{Pre-Submission Sequence and Evidence Package}

\draftnote{Order the pre-submission steps and name the evidence each one
carries, drawing the credibility evidence from
\mvfile{funding/pdac-funding-applications/final-apply/sections/sec-05-trial-evidence.tex}:
the ASME V\&V 40 and ICH M15 aligned credibility score of 81.9 across 55
verification notebook tests. Carry the limitation each source stated in the same
row as its result.}
